%% notes-tex/030-scaling-study/sections/3-readout.tex

\section{The readout\secstatus{todo}}
\label{sec:readout}

The readout (\texttt{scaling\_readout\_030.py}) finds every run of a sweep by its
config tag, stitches chained continuations by the checkpoint each link resumed
from, and keeps one row per epoch of every point with the validation loss in
label units, pooled and per screen, the Pearson beside it, and the compute keys
(Figure~\ref{fig:readout}). A point is one (size, subset, seed, learning rate).

\tcfig{figures/030-scaling-readout.pdf}{%
The readout from W\&B rows to fitted exponents. The five boxes are the stages;
the purple box is the 025 paired sweep, which enters the same table and leaves as
the gap; the red box is the replicate floor, computed before any fit and drawn
beside every fitted $E$. The dashed boxes are the four joints repaired on
2026-10-12, each of which had been reproduced on a real run or a real compose
before it was changed (Section~\ref{sec:support}).
}{fig:readout}
\src{notes/assets/drawio/030-scaling-study.gen.py}

\notesec{The loss that is fitted}
The fit target is the mean squared error on the gene-interaction score after the
inverse transform, \texttt{val/gene\_interaction/MSE}, and per screen the
per-token keys \texttt{val/token/[screen]/gene\_interaction/MSE}. Not the
normalized loss, whose constants differ per arm, and not Pearson, which plateaus
and carries a floor of its own: the two screens' validation triples differ in
mean by about one standard deviation (Kuzmin 2018 mean $-0.046$, sd 0.049,
$n = 10{,}067$; Kuzmin 2020 mean $+0.004$, sd 0.062, $n = 28{,}935$, measured on
the Q validation entries), so predicting only the per-screen mean scores a pooled
Pearson of 0.34 and every pooled Pearson carries that component. Squared error
does not, but the two screens still have different replicate floors, so the
curves are fitted per screen and the pooled curve is read as their mixture. The
per-token MSE is computed over validation entry rows (39,002 rows of the 37,705
triples), the pooled MSE over the policy-reduced one value per triple; the
populations differ by the triples measured more than once in one screen.

\notesec{Two reads per point}
The selected read is the epoch of lowest pooled validation MSE. For the subsets
that converge inside their budget (240 epochs of 1/64, 120 of 1/16) it is the
converged loss; for the full pool at ten epochs it is the last epoch or close to
it. The matched read interpolates each point's epoch curve at one records-seen
budget $B$ shared by every point. $B$ defaults to the smallest total any point
reached, which the budgets make the 1/64 total of 5.12M records: the full pool at
ten epochs sees 10.46M, 1/4 at 40 epochs 9.78M, 1/16 at 120 epochs 8.26M, so a
budget set at the full-data total would leave every subset without the read,
which is how the first version of the readout was written. At $B = 5.12$M the
full point is read near its fifth epoch and 1/4 near its twenty-first. The
matched read is loss against unique data at a fixed training budget, the
data-constrained framing in which repeated epochs of a small subset are set
against single passes over a large one.

\notesec{Fits}
On each read and for each target (pooled, 2018, 2020), two forms on each axis:
the pure power law $L = A x^{-k}$ and the floored form $L = E + A x^{-k}$, where
$x$ is $D$ in records, $D$ in entry rows (030 only) at the 180x8 column, and $N$
at the full-pool column. The fit is least squares in log space from a grid of
starting exponents, with a Huber residual at $\delta = 10^{-3}$ as Hoffmann
et~al.\ fit, and $E$ bounded between zero and the smallest observed loss. For
every fit the largest $x$ is held out, the form is refit on the rest, and its
prediction at the held-out point is reported beside the observation; the floored
form has three parameters, so this test needs four distinct $x$ in sample, and
on three it is an exact fit that predicts nothing. Three hundred bootstrap
resamples over runs give the 5 to 95 percent interval of every exponent. The
interval, not the point, is what an extrapolation inherits.

\notesec{The floor before any fit}
Under a squared loss the best predictor of a single noisy measurement has
expected error equal to the replicate variance, so the released and re-measured
records of the screens bound $E$ from below in the loss's own units before a fit
exists (Table~\ref{tab:floors}). The 2020 reliability is measured on this build's
11,013 same-strain re-measured triples; 2018 has 61 such pairs, so it is bounded
by the published Dango replicate correlation below and by the 483 cross-year
pairs above. On 2020 the floor is 56\% of the variance of the score: no predictor
removes more than 44\% of the loss of predicting the screen mean. A fitted $E$
below the floor is a wrong fit; one far above it says the model family, not the
data, is the limit.

%% GENERATED FILE -- do not hand-edit.
%% SOURCE: experiments/030-solid-growth-multi/scripts/scaling_grid_table_030.py, reading triple_noise_ceiling_030.json (noise_model sd_of_scores, pair_classes pearson)

\begin{table}[htbp]
\centering\footnotesize
\caption{The replicate floor per screen in label units. Under a squared loss the best predictor of a single noisy measurement has expected error $E = (1 - r)\,\mathrm{var}(\tau)$ where $r$ is the replicate reliability, the correlation of two independent measurements of the same triple. The variance column is the loss of predicting the screen mean, the top of the range a curve can move through, so the floor's share of it is the fraction of the loss no model can remove. The 2020 reliability is measured on this build's same-strain re-measured pairs; 2018 has 61 such pairs, so it is bounded by the published Dango replicate correlation below and the cross-year 2018 vs 2020 pairs above.}
\label{tab:floors}
\begin{tabular}{llrrrrr}
\toprule
screen & reliability from & $r$ & sd($\tau$) & var($\tau$) & $E$ & $E$/var \\
\midrule
Kuzmin 2018 & Dango 0.59 & 0.59 & 0.0535 & 2.86e-3 & 1.17e-3 & 0.41 \\
Kuzmin 2018 & cross-year 2018 vs 2020 pairs & 0.89 & 0.0535 & 2.86e-3 & 0.31e-3 & 0.11 \\
Kuzmin 2020 & same-strain re-measured pairs & 0.44 & 0.0633 & 4.01e-3 & 2.25e-3 & 0.56 \\
\bottomrule
\end{tabular}
\end{table}


\notesec{The paired gap}
The 025 sweep passes through the same table and the same two reads at the same
$B$. The gap at each subset and read is the 030 pooled MSE minus the 025 pooled
MSE, each a mean over seeds with the seed standard deviation where there are
three. It is pooled only: the 025 trainer logs no per-screen key, and its
validation triples are the same genotypes so a per-screen 025 loss would need
its predictions dumped and scored offline.
